% Source: https://artofproblemsolving.com/community/c3h2013942_random_geometry_theorem
% Stable identity: aops-topic-2013942-q01
The angles at $B$ and $C$ are acute (they add to $90^\circ$), so the foot $D$ lies strictly between $B$ and $C$.

\begin{asy}
usepackage("amsmath");
size(5.5cm);
defaultpen(fontsize(10pt));
pair B=(0,0), C=(6.5,0); real t=1.8;
pair A=(t,sqrt(t*(6.5-t)));
pair D=(t,0);
draw(A--B--C--cycle,linewidth(0.9));
draw(A--D,red+linewidth(0.9));
draw(D+(0.25,0)--D+(0.25,0.25)--D+(0,0.25));
pair ua=unit(B-A), ub=unit(C-A);
draw(A+0.3*ua--A+0.3*ua+0.3*ub--A+0.3*ub);
dot(A); dot(B); dot(C); dot(D);
label("$A$",A,N); label("$B$",B,SW); label("$C$",C,SE); label("$D$",D,S);
label("\mbox{}",(min(currentpicture,user=true)+max(currentpicture,user=true))/2,black); // keeps the text after the figure black when LaTeX typesets the labels inline
\end{asy}

Compare the right triangles $DBA$ and $DAC$, both with a right angle at $D$. In triangle $ABD$, $\angle DAB=90^\circ-\angle B$. In triangle $ABC$, $\angle C=90^\circ-\angle B$. Hence
\[
\angle DAB=\angle C=\angle DCA,
\]
and similarly $\angle DBA=\angle DAC$. So triangles $DBA$ and $DAC$ are similar, with $D\leftrightarrow D$, $B\leftrightarrow A$, $A\leftrightarrow C$. Corresponding sides are proportional:
\[
\dfrac{BD}{AD}=\dfrac{AD}{CD},
\]
and cross-multiplying gives $AD^2=BD\cdot CD$.
